\documentclass{article}

\usepackage{fancyhdr}
\usepackage{extramarks}

%amsmath
\usepackage{amsmath}
\usepackage{amsthm}
\usepackage{amsfonts}
\usepackage{mathtools}
\usepackage{amssymb}

%Enumerating items and adding columns
\usepackage{enumitem}
\usepackage{multicol}

%TikZ picture
\usepackage{tikz}
\usepackage{scalerel}
\usepackage{pict2e}
\usepackage{tkz-euclide}
\usetikzlibrary{calc}
\usetikzlibrary{patterns,arrows.meta}
\usetikzlibrary{shadows}
\usetikzlibrary{external}
\usetikzlibrary{automata,positioning}

%pgfplots
\usepackage{pgfplots}
\pgfplotsset{compat=newest}
\usepgfplotslibrary{statistics}
\usepgfplotslibrary{fillbetween}

%colours
\usepackage{xcolor}

%
% Basic Document Settings
%

\topmargin=-0.45in
\evensidemargin=0in
\oddsidemargin=0in
\textwidth=6.5in
\textheight=9.0in
\headsep=0.25in

\linespread{1.1}

\pagestyle{fancy}
\lhead{\hmwkAuthorName}
\chead{\hmwkChapter . \hmwkChapterTitle \hmwkClassTime: \hmwkTitle}
\rhead{\firstxmark}
\lfoot{\lastxmark}
\cfoot{\thepage}

\renewcommand\headrulewidth{0.4pt}
\renewcommand\footrulewidth{0.4pt}

\setlength\parindent{0pt}

%
% Create Problem Sections
%

\newcommand{\enterProblemHeader}[1]{
    \nobreak\extramarks{}{Problem \arabic{#1} continued on next page\ldots}\nobreak{}
    \nobreak\extramarks{Problem \arabic{#1} (continued)}{Problem \arabic{#1} continued on next page\ldots}\nobreak{}
}

\newcommand{\exitProblemHeader}[1]{
    \nobreak\extramarks{Problem \arabic{#1} (continued)}{Problem \arabic{#1} continued on next page\ldots}\nobreak{}
    \stepcounter{#1}
    \nobreak\extramarks{Problem \arabic{#1}}{}\nobreak{}
}

\setcounter{secnumdepth}{0}
\newcounter{partCounter}
\newcounter{homeworkProblemCounter}
\setcounter{homeworkProblemCounter}{1}
\nobreak\extramarks{Problem \arabic{homeworkProblemCounter}}{}\nobreak{}

%
% Homework Problem Environment
%
% This environment takes an optional argument. When given, it will adjust the
% problem counter. This is useful for when the problems given for your
% assignment aren't sequential. See the last 3 problems of this template for an
% example.
%
\newenvironment{homeworkProblem}[1][-1]{
    \ifnum#1>0
        \setcounter{homeworkProblemCounter}{#1}
    \fi
    \section{Problem \arabic{homeworkProblemCounter}}
    \setcounter{partCounter}{1}
    \enterProblemHeader{homeworkProblemCounter}
}{
    \exitProblemHeader{homeworkProblemCounter}
}

%
% Homework Details
%   - Title
%   - Due date
%   - Class
%   - Section/Time
%   - Instructor
%   - Author
%

\newcommand{\hmwkTitle}{Exercises}
\newcommand{\hmwkChapter}{CHAPTER 3}
\newcommand{\hmwkDate}{April 3, 2026}
\newcommand{\hmwkClassTime}{}
\newcommand{\hmwkChapterTitle}{Sets}
\newcommand{\hmwkAuthorName}{\textbf{Christian Jelo R. Artoza}}

%
% Title Page
%

\title{
    \vspace{2in}
    \textmd{\textbf{\hmwkChapter:\ \hmwkTitle}}\\
    \normalsize\vspace{0.1in}\small{\hmwkDate}\\
    \vspace{0.1in}\large{\textit{\hmwkChapterTitle\ \hmwkClassTime}}
    \vspace{3in}
}

\author{\hmwkAuthorName}
\date{}

\renewcommand{\part}[1]{\textbf{\large Part \Alph{partCounter}}\stepcounter{partCounter}\\}

%
% Various Helper Commands
%

% qedsymbol that is always flushed to the right
\newcommand{\qedright}{
    \begin{flushright}
        \qedsymbol
    \end{flushright}
}

% For propositions
\newtheorem{prop}{Proposition}

% Useful for algorithms
\newcommand{\alg}[1]{\textsc{\bfseries \footnotesize #1}}

% For derivatives
\newcommand{\deriv}[1]{\frac{\mathrm{d}}{\mathrm{d}x} (#1)}

% For partial derivatives
\newcommand{\pderiv}[2]{\frac{\partial}{\partial #1} (#2)}

% Integral dx
\newcommand{\dx}{\mathrm{d}x}

% Alias for the Solution section header
\newcommand{\solution}{\textbf{\large Solution}}

% Probability commands: Expectation, Variance, Covariance, Bias
\newcommand{\E}{\mathrm{E}}
\newcommand{\Var}{\mathrm{Var}}
\newcommand{\Cov}{\mathrm{Cov}}
\newcommand{\Bias}{\mathrm{Bias}}

\begin{document}

\maketitle

\pagebreak

%Plot Settings
\pgfplotsset{
    standard/.style={
    axis line style = thick,
    axis equal,
    trig format=rad,
    enlargelimits,
    axis x line=middle,
    axis y line=middle,
    axis line style={latex-latex},
    enlarge x limits=0.15,
    enlarge y limits=0.15,
    every axis x label/.style={at={(current axis.right of origin)},anchor=north west},
    every axis y label/.style={at={(current axis.above origin)},anchor=south east}
    }
}

\begin{homeworkProblem}
    \textit{(Exercise 1.22.)} Given an example of 100 numbers from $\{1,\,2,\,3,\dots,\,200\}$
    where not one of your numbers divides another. This proves that Proposition 1.11 is optimal.
    \vspace{1.5mm}

    \textbf{Solution.} Start by choosing all the larger odd numbers $m$ such that $3m > 200$. There are
    67 such numbers; these are:
    \begin{align*}
        67&,69,71,73,75,77,79,81,83,85,87,89,91,93,95,97,99,101,103,105 \\
        107&,109,111,113,115,117,119,121,123,125,127,129,131,133,135,137 \\
        139&,141,143,145,147,149,151,153,155,157,159,161,163,165,167,169 \\
        171&,173,175,177,179,181,183,185,187,189,191,193,195,197,199
    \end{align*}
    For each $m$, the multiple $2m$ is not in the selection since all are odd numbers. Furthermore,
    since the multiple $3m > 200$, it follows that $3m$ and the larger multiples are not included.
    Thus, no number in the selection divides another.
    \vspace{1.5mm}

    Next, we select the numbers $2m$ where $m$ is an odd number before $67$ such that that $3m \geq 67$.
    There will be 22 such numbers, which are the following:
    \begin{align*}
        46&,50,54,58,62,66,70,74,78,82,86,90,\\
        94&,98,102,106,110,114,118,122,126,130
    \end{align*}
    Since these are even numbers, none of these numbers divide those in the first selection. Furthermore,
    since $3m \geq 67$ for each $2m$ in this selection, it follows that $3\cdot 2m \geq 134$. Thus,
    none of the larger multiples of each number are included in the selection. Hence, none divides
    another.
    \vspace{1.5mm}

    We continue selecting in this fashion: selecting the numbers $2^k \cdot m$ where $m$ is an odd
    number before the odd number $p$ corresponding to the smallest number, $2^\ell \cdot p$
    for some $\ell$, in the previous selection; with the property that $3m \geq p$. Doing so
    gives the following 100 numbers.
    \begin{align*}
        16&,24,40,56,36,44,52,60,68,76,84,46,50,54,58,62,66,70,74,78,82,86,90,94, \\
        98&,102,106,110,114,118,122,126,130,67,69,71,73,75,77,79,81,83,85,87,89,91,\\
        93&,95,97,99,101,103,105,107,109,111,113,115,117,119,121,123,125,127,129,\\
        131&,133,135,137,139,141,143,145,147,149,151,153,155,157,159,161,163,165,\\
        167&,169,171,173,175,177,179,181,183,185,187,189,191,193,195,197,199
    \end{align*}
    By construction, no number in the selection divides another. \qedright
\end{homeworkProblem}

\begin{homeworkProblem}
    \textit{(Exercise 1.23.)} Suppose you deal a pile of cards, face down, from a shuffled deck of
    cards (this is a standard 52-card deck, where each card is one of 4 suits and one of 13 ranks).
    How many must you deal out until you are guaranteed...
    \begin{enumerate}
        \item five of the same suit? That is, a flush.
        \item two of the same rank? That is, a two-of-a-kind.
        \item three of the same rank? That is, a three-of-a-kind.
        \item four of the same rank? That is, a four-of-a-kind.
        \item two of one rank and three of another? That is, a full house.
    \end{enumerate}
    \textbf{Solution.}

    1. There are four suits in a deck of cards (Clubs, Spades, Hearts, Diamonds). So, by the Pigeonhole
    princeiple, five of the same suit is guaranteed when you deal:
    \begin{align*}
        (4)(4) + 1 = 17\text{ cards}
    \end{align*}
    2. There are thirteen ranks in a deck of cards. So, in order to guarantee a two-of-a-kind, you have
    to deal:
    \begin{align*}
        (13)(1) + 1 = 14\text{ cards}
    \end{align*}
    3. A three-of-a-kind is guaranteed by dealing:
    \begin{align*}
        (13)(2) + 1 = 27\text{ cards}
    \end{align*}
    4. A four-of-a-kind is guaranteed by dealing:
    \begin{align*}
        (13)(3) + 1 = 40\text{ cards}
    \end{align*}
    5. Pair up each rank to another rank. For instance, consider pairing Ace and Two, Three and Four,
    Five and Six, Seven and Eight, Nine and Ten, Queen and Jack, but leave Kings without a pair.
    We'll create a box for each pair and a box for the Kings.
    \vspace{5mm}

    \fbox{
        \begin{minipage}[c][8ex]{1.5cm}
            \begin{center}
                Ace \& 2
            \end{center}
        \end{minipage}
    }\quad
    \fbox{
        \begin{minipage}[c][8ex]{1.5cm}
            \begin{center}
                3 \& 4
            \end{center}
        \end{minipage}
    }\quad
    \fbox{
        \begin{minipage}[c][8ex]{1.5cm}
            \begin{center}
                5 \& 6
            \end{center}
        \end{minipage}
    }\quad
    \fbox{
        \begin{minipage}[c][8ex]{1.5cm}
            \begin{center}
                7 \& 8
            \end{center}
        \end{minipage}
    }\quad
    \fbox{
        \begin{minipage}[c][8ex]{1.5cm}
            \begin{center}
                9 \& 10
            \end{center}
        \end{minipage}
    }\quad
    \fbox{
        \begin{minipage}[c][8ex]{1.5cm}
            \begin{center}
                Jack \& Queen
            \end{center}
        \end{minipage}
    }\quad
    \fbox{
        \begin{minipage}[c][8ex]{1.5cm}
            \begin{center}
                King
            \end{center}
        \end{minipage}
    }
    \vspace{5mm}

    It is likely that we could deal enough cards to get three cards of the same rank, and deal more
    cards enough that we get a two-of-a-kind whose rank is paired with that of the three-of-a-kind,
    giving us a full house. But, it is also possible to get a four-of-a-kind before getting the
    two-of-a-kind. So, in order to guarantee a full house, one box must contain 6 cards. Thus, by the
    Pigeonhole principle, we need to deal
    \begin{align*}
        ((6)(5)+1)+4 = 35\text{ cards. }
    \end{align*}
    \qedright
\end{homeworkProblem}

\begin{homeworkProblem}
    \textit{(Exercise 1.24.)} Prove that any set of seven integers contains a pair whose sum or
    difference is divisible by 10.
    \begin{proof}
        First, note that divisibility of a sum $n+m$, or difference $n-m$, by 10 depends on the
        ones digit of $n$ and $m$. In particular, if the sum of the ones digit of two numbers is equal
        to 10, then the sum of those numbers is divisible by 10. Furthermore, if the difference of
        the ones digit of two numbers is equal to 0, then the difference of those numbers is divisible
        by 10. Using this fact, we create a box for each pair of number that add up to 10; while leaving
        the numbers 0 and 5 alone in a box.
        \begin{center}
            \fbox{
            \begin{minipage}[c][8ex]{1.5cm}
                \begin{center}
                    0
                \end{center}
            \end{minipage}
            }\quad
            \fbox{
                \begin{minipage}[c][8ex]{1.5cm}
                    \begin{center}
                        1 or 9
                    \end{center}
                \end{minipage}
            }\quad
            \fbox{
                \begin{minipage}[c][8ex]{1.5cm}
                    \begin{center}
                        2 or 8
                    \end{center}
                \end{minipage}
            }\quad
            \fbox{
                \begin{minipage}[c][8ex]{1.5cm}
                    \begin{center}
                        3 or 7
                    \end{center}
                \end{minipage}
            }\quad
            \fbox{
                \begin{minipage}[c][8ex]{1.5cm}
                    \begin{center}
                        4 or 6
                    \end{center}
                \end{minipage}
            }\quad
            \fbox{
                \begin{minipage}[c][8ex]{1.5cm}
                    \begin{center}
                        5
                    \end{center}
                \end{minipage}
            }
        \end{center}
        Then, given a set of seven integers, we put each integer into the boxes according to their ones
        digit. For instance, if the integer has 9 as its ones digit, then we put it into the box
        corresponding to 1 or 9. With 7 integers to put into 6 boxes, it follows by the Pigeonholse
        principle that at least one box contains at least two integers $n,m$. If $n$ and $m$ are
        contained in the "0" or "5" box, then both their sum and difference are divisible by 10. Suppose
        $n$ and $m$ are contained in the other boxes. By construction of the boxes, the sum of their ones
        digits must be equal to 10, if their ones digits are not the same. Hence, $n+m$ is divisible by
        10. But, if the ones digits are the same, then the difference of their once digits is equal to 0.
        Hence, $n-m$ is divisible by 10.
    \end{proof}
\end{homeworkProblem}

\begin{homeworkProblem}
    \textit{(Exercise 1.25.)} Read the \textit{Introduction to Ramsey theory} following this
    chapter. Then, let $r(n,m)$ be the smallest value $N$ for which every red/blue coloring
    of $K_N$ contains either a red $K_n$ or a blue $K_m$. Prove that $r(n,2)=n$
    \begin{proof}
    We want to prove that $n$ is the smallest value such that every red/blue coloring
    of $K_n$ contains either a red $K_n$ or a blue $K_2$. Clearly it can not be smaller than $n$
    since the coloring of $K_{n-1}$ where all edges are colored red neither contains a red $K_n$ nor a
    blue $K_2$. Furthermore, $K_n$ is either colored all red, or not. If it is not colored all red,
    then it contains a blue $K_2$. Therefore, $r(n,2) = n$, as desired.
    \end{proof} 
\end{homeworkProblem}

\begin{homeworkProblem}
    \textit{(Exercise 1.26.)} Determine the U.S. population at the time that you are reading
    this.
    \begin{enumerate}
        \item Does the pigeonhole principle guarantee that 1 million U.S. residents all have
        the same birthday?
        \item If the principle does not guarantee this, how many people are needed for until that
        milestone is reached? If the USA grows by 2 million people per year, in what year will
        this occur?
    \end{enumerate}
    \textbf{Solution.} The current population of the US as of 17:43 UTC March 30, 2026 is 342,407,233.
    \begin{enumerate}
        \item It is guaranteed that 1 million U.S. residents all have the same birthday if the total
        population is:
        \begin{align*}
            (366)(999,999)+1=365,999,635
        \end{align*}
        Comparing this with the current population, it follows that the pigeonhole principle does not
        guarantee that 1 million residents have the same birthday.

        \item We would need:
        \begin{align*}
            365,999,635 - 342,407,233 = 23,592,402
        \end{align*}
        more U.S. residents. And, if the USA grows by 2 million people per year, then the milestone
        will be reached after
        \begin{align*}
            \frac{23,592,402}{2,000,000} \approx 12
        \end{align*}
        years.\qedright
    \end{enumerate}
\end{homeworkProblem}

\end{document}